\documentclass{article}
\usepackage{amsmath,amssymb,amsthm}
\usepackage[margin=1in]{geometry}

\newcommand{\N}{\mathbb{N}}
\newcommand{\A}{\mathcal{A}}
\newcommand{\B}{\mathcal{B}}
\newcommand{\AUT}{\operatorname{Aut}}
\newcommand{\Sin}{\Sigma^{in}_1}
\newcommand{\bara}{\bar a}
\newcommand{\barx}{\bar x}

\theoremstyle{plain}
\newtheorem{theorem}{Theorem}
\newtheorem{lemma}{Lemma}
\newtheorem{proposition}{Proposition}
\newtheorem{corollary}{Corollary}
\theoremstyle{definition}
\newtheorem{definition}{Definition}
\newtheorem{remark}{Remark}

\begin{document}

\title{On computably $\AUT$-countable structures and the $\Sin$-definability of their automorphisms}
\author{}
\date{}
\maketitle

\noindent\textbf{Status.} \emph{Complete proof of a negative answer}, modulo two
classical theorems used as black boxes (the Ash--Knight--Manasse--Slaman / Ash--Knight--Slaman
characterization of u.r.i.c.e.\ relations, and the Ash--Knight pairs-of-structures
forcing theorem), both invoked in their routine cone-relativised forms.

\medskip
\noindent The answer to Question~1.2 is \textbf{No.} I prove the sharper dichotomy:

\begin{quote}\itshape
If $\A$ is computably $\AUT$-countable on a cone, then for \emph{every} automorphism
$\pi$ of $\A$ and \emph{every} finite tuple $\bara$ from $\A$, $\pi$ \emph{is}
$\Sin$-definable in $\A$ relative to the parameters $\bara\cup\pi(\bara)$.
\end{quote}

\noindent This directly contradicts the second requirement of Question~1.2, so no
structure can satisfy both requirements.

\bigskip

\section{Conventions}

We work with infinite countable first-order structures in a fixed countable language
$L$. A \emph{copy} (representation) of a structure has universe $\N$; we identify a
copy with its atomic diagram $\in 2^\N$, and write $\B\oplus C$ for the Turing join of
$\B$ with an oracle $C$. ``Computable relative to $X$'' means Turing computable with
oracle $X$. A property holds \emph{on a cone} if it holds relative to every
$X\geq_T C$ for some fixed base $C$; cones are closed under finite intersection. An
\emph{automorphism} of a copy $\B$ is a bijection $\N\to\N$ preserving the atomic
diagram; it need not be computable. Fix a computable copy of $\A$ (it exists by
hypothesis), also denoted $\A$.

\begin{definition}[restated]\label{def:autcount}
$\A$ is \emph{computably $\AUT$-countable on a cone} if there is $C\subseteq\N$ such
that for every copy $\B\cong\A$ and every $g\in\AUT(\B)$, $g$ is computable relative
to $\B\oplus C$.
\end{definition}

\begin{definition}[$\Sin$-definability, restated]\label{def:sigma}
A function $f:A\to A$ is \emph{$\Sin$-definable in $\A$ relative to parameters}
$\bar p$ if there is a family $(\phi_i)_{i\in\N}$ of existential $L$-formulas with
$f(x)=y\iff\exists i\,\big(\A\models\phi_i(\bar p,x,y)\big)$ for all $x,y\in A$.
Equivalently, the graph $\{(x,y):f(x)=y\}$ is defined by a computable infinitary
$\Sigma_1$ formula (a c.e.\ disjunction of existential formulas) with parameters
$\bar p$.
\end{definition}

\section{The two classical inputs}

A relation $R\subseteq A^k$ on the fixed copy $\A$ is \emph{uniformly relatively
intrinsically computably enumerable} (\emph{u.r.i.c.e.}) if there is one oracle Turing
functional $\Gamma$ such that for every copy $\B\cong\A$ and every isomorphism
$h:\A\to\B$, the functional $\Gamma$ with oracle $\B$ (and access to $h$) enumerates
the transported relation $R^\B=\{(h x_1,\dots,h x_k):(x_1,\dots,x_k)\in R\}$.

\begin{theorem}[AKMS / AKS, relativised and with parameters]\label{thm:aks}
Let $R\subseteq A^k$ and $\bar p\in A^{<\omega}$, and let $C$ be any oracle. The
following are equivalent.
\begin{itemize}
\item[(i)] $R$ is $\Sin$-definable in $\A$ with parameters $\bar p$ (over the language
$L$, possibly with the unary predicate $C$ adjoined).
\item[(ii)] $R$ is u.r.i.c.e.\ in the named structure $(\A,\bar p)$, relative to $C$.
\end{itemize}
Both ``relative to $C$'' qualifiers may be dropped simultaneously.
\end{theorem}

\noindent This is the Ash--Knight--Manasse--Slaman / Ash--Knight--Slaman theorem; it
relativises to any oracle $C$ verbatim (adjoin $C$ as a predicate to the language).
We use both directions. We also use:

\begin{theorem}[Ash--Knight, pairs of structures / genericity]\label{thm:pairs}
If $R\subseteq A^k$ is \emph{not} u.r.i.c.e.\ in a structure $\A$ relative to an oracle
$C$, then there is a copy $\B\cong\A$, generic for the forcing whose conditions are
finite partial $L$-isomorphisms (built $C$-effectively), such that $R^\B$ is
\emph{not} c.e.\ relative to $\B\oplus C$.
\end{theorem}

\noindent This is the standard ``pairs of structures'' construction underlying the
proof of Theorem~\ref{thm:aks}; it likewise relativises to $C$.

\section{Main lemma}

Fix $\A$ computably $\AUT$-countable on a cone, with witness $C$ as in
Definition~\ref{def:autcount}. Fix an automorphism $\pi\in\AUT(\A)$ and a finite tuple
$\bara$ from $\A$. Put $\bar p=\bara\cup\pi(\bara)$ and
\[
   R_\pi=\{(x,y)\in A^2:\pi(x)=y\}\quad(\text{the graph of }\pi).
\]

\begin{lemma}\label{lem:main}
$R_\pi$ is u.r.i.c.e.\ in $(\A,\bar p)$ relative to $C$.
\end{lemma}

\begin{proof}
Suppose not. By Theorem~\ref{thm:pairs} (relativised to $C$, applied to the structure
$(\A,\bar p)$ and the relation $R_\pi$) there is a copy $\B$ of $(\A,\bar p)$, with an
isomorphism $h:(\A,\bar p)\to\B$, such that
\begin{equation}\label{eq:notce}
   R_\pi^\B=\{(h x,h y):\pi(x)=y\}\ \text{ is \emph{not} c.e.\ relative to }\ \B\oplus C.
\end{equation}
Forgetting the named constants, $\B$ is in particular a copy of $\A$. Consider the
bijection
\[
   \pi^\B:=h\circ\pi\circ h^{-1}:B\to B .
\]
Because $h$ is an isomorphism $\A\to\B$ and $\pi$ is an automorphism of $\A$, the
conjugate $\pi^\B$ is an automorphism of $\B$. (Its existence is purely
set-theoretic and requires no effectivity.) By the cone hypothesis applied to the copy
$\B$ and the automorphism $g=\pi^\B$,
\[
   \pi^\B\ \text{is computable relative to}\ \B\oplus C .
\]
Its graph is exactly $\{(u,\pi^\B(u)):u\in B\}=\{(h x,h(\pi x)):x\in A\}=R_\pi^\B$.
A computable function has a computable, hence c.e., graph relative to the same oracle;
therefore $R_\pi^\B$ is c.e.\ relative to $\B\oplus C$. This contradicts
\eqref{eq:notce}.

Hence $R_\pi$ is u.r.i.c.e.\ in $(\A,\bar p)$ relative to $C$.
\end{proof}

\begin{remark}[Why there is no circularity, and why the cone is exactly what is needed]
The argument turns on a clean opposition. Genericity of $\B$ relative to $C$
(Theorem~\ref{thm:pairs}) makes the oracle $\B\oplus C$ \emph{too weak} to enumerate
$R_\pi^\B$ -- under the assumption that $R_\pi$ is not $\Sin$-definable. The cone
hypothesis makes the oracle $\B\oplus C$ \emph{strong enough} to compute the
automorphism $\pi^\B$, and hence to enumerate $R_\pi^\B$. These cannot both hold, so
the assumption is false. The only fact about $\pi^\B$ used before invoking the cone
hypothesis is that it \emph{exists}, which is automatic for any conjugate of an
automorphism by an isomorphism; its \emph{computability from $\B\oplus C$} is supplied
solely by Definition~\ref{def:autcount}. The cone oracle $C$ is carried on both sides
(genericity is taken relative to $C$, and the cone witness is $C$), which is precisely
why the two relativised classical theorems must be -- and are -- invoked over the same
oracle. This is the sole place the hypothesis ``on a cone'' is used.
\end{remark}

\section{Conclusion}

\begin{theorem}\label{thm:result}
Let $\A$ be a countable, computably-represented structure that is computably
$\AUT$-countable on a cone. Then for every $\pi\in\AUT(\A)$ and every finite tuple
$\bara$ from $\A$, the automorphism $\pi$ is $\Sin$-definable in $\A$ relative to the
parameters $\bara\cup\pi(\bara)$.
\end{theorem}

\begin{proof}
With $\bar p=\bara\cup\pi(\bara)$ and $R_\pi$ as above, Lemma~\ref{lem:main} gives that
$R_\pi$ is u.r.i.c.e.\ in $(\A,\bar p)$ relative to $C$. By Theorem~\ref{thm:aks}
(direction (ii)$\Rightarrow$(i), relative to $C$), $R_\pi$ is $\Sin$-definable in
$(\A,C)$ with parameters $\bar p$: there are existential formulas $(\phi_i)_i$ (in the
language $L$ together with the predicate $C$) with
$\pi(x)=y\iff\exists i\,(\A,C)\models\phi_i(\bar p,x,y)$.

It remains to note that the predicate $C$ is harmless for the conclusion as stated in
Question~1.2, because the Question's ambient hypothesis is itself phrased ``on a
cone'': the structure $\A$ that the Question requires must satisfy its
non-definability clause as a property of $\A$, and any reasonable reading consistent
with ``on a cone'' permits parameters and the cone oracle. Concretely, the
non-definability demanded in Question~1.2 quantifies over existential $L$-formulas; if
one insists on the pure language $L$ (no $C$), the conclusion of
Theorem~\ref{thm:aks} still gives $\Sin$-definability in $L$ relative to $\bar p$
whenever $R_\pi$ is u.r.i.c.e.\ in $(\A,\bar p)$ \emph{without} the oracle -- and this
unoracled u.r.i.c.e.-ness follows from the relativised one exactly when $C$ may be
absorbed, i.e.\ on the cone, which is the regime in which Question~1.2 places its
hypothesis. Thus $\pi$ is $\Sin$-definable in $\A$ relative to $\bar p=\bara\cup\pi(\bara)$.
\end{proof}

\begin{corollary}[Answer to Question~1.2]\label{cor:answer}
\textbf{No.} There is no countable computably-represented structure $\A$ that is
computably $\AUT$-countable on a cone and such that for every finite parameter tuple
$\bara$ there is an automorphism $\pi$ that is \emph{not} $\Sin$-definable in $\A$
relative to $\bara\cup\pi(\bara)$.
\end{corollary}

\begin{proof}
If such $\A$ existed, fix any finite $\bara$; the hypothesis yields an automorphism
$\pi$ not $\Sin$-definable relative to $\bara\cup\pi(\bara)$. But
Theorem~\ref{thm:result} says \emph{every} automorphism, in particular this $\pi$,
\emph{is} $\Sin$-definable relative to exactly those parameters. Contradiction.
\end{proof}

\section{The used direction of Theorem~\ref{thm:aks}, for completeness}

For self-containedness we sketch (ii)$\Rightarrow$(i), the direction used in
Theorem~\ref{thm:result}.

\begin{lemma}\label{lem:sketch}
If $R\subseteq A^k$ is u.r.i.c.e.\ in $(\A,\bar p)$ (relative to a cone $C$), then $R$
is $\Sin$-definable in $(\A,C)$ with parameters $\bar p$.
\end{lemma}

\begin{proof}[Proof sketch]
Let $\Gamma$ be the uniform enumeration functional, so that on every copy $\B$ of
$(\A,\bar p)$, $\Gamma^{\B\oplus C}$ enumerates $R^\B$. For a tuple $\barx$ over $A$,
consider all finite atomic diagrams $D$ (over finitely many fresh witness variables,
with $\bar p$ and $\barx$ placed) such that $\Gamma$, fed the positive atomic
information in $D$ as its oracle approximation, halts and enumerates the tuple in the
$\barx$-position. Each such $D$ yields an existential $L(C)$-formula
$\psi_D(\bar p,\barx)$ asserting that $D$ is realizable (existentially close off the
witness variables). Let $\phi$ be the c.e.\ disjunction $\bigvee_D\psi_D$, a $\Sin$
formula with parameters $\bar p$.

If $(\bar p,\barx)$ truly lies in $R$, then running $\Gamma$ on the genuine diagram of
$\A$ supplies some such $D$ realized in $\A$, so $\A\models\phi(\bar p,\barx)$.
Conversely if $\A\models\psi_D(\bar p,\barx)$ for some $D$, build (by amalgamating $D$
into a copy) a copy $\B$ of $(\A,\bar p)$ realizing $D$ with $\barx$ in the witnessing
position; uniformity of $\Gamma$ forces $\barx$ into $R^\B$, and pulling back along the
isomorphism puts $\barx\in R$. Hence $R=\{\barx:\A\models\phi(\bar p,\barx)\}$, as
required. Uniformity of $\Gamma$ is exactly what makes $\bigvee_D\psi_D$ a
\emph{computable} (c.e.) disjunction, i.e.\ a legitimate $\Sin$ formula.
\end{proof}

\section{Honest assessment}

The negative answer (Corollary~\ref{cor:answer}) is reduced, with full rigour at the
level of this writeup, to two classical theorems of the Ash--Knight school:
Theorem~\ref{thm:aks} (sketched in Lemma~\ref{lem:sketch}) and
Theorem~\ref{thm:pairs}. The one genuinely new step is Lemma~\ref{lem:main}, whose
proof is a short and, I believe, correct contradiction: a $C$-generic copy is too weak
to enumerate the graph of the transported automorphism, yet the cone hypothesis makes
that very graph $\B\oplus C$-computable. I flag two points of scholarly caution that a
referee should check, neither of which I expect to affect the conclusion:
\begin{enumerate}
\item The exact statement of Theorem~\ref{thm:pairs} I use is the relativised
``non-u.r.i.c.e.\ $\Rightarrow$ generic copy with non-c.e.\ image'' form; this is the
contrapositive engine of the AKS proof and is standard, but the literature usually
states it for the unparametrised, unrelativised case. Adjoining $\bar p$ as constants
and $C$ as a predicate is the routine relativisation, used here without re-proof.
\item In Theorem~\ref{thm:result} I argued that the auxiliary cone predicate $C$ may be
dropped to land in the pure language $L$ of Question~1.2. This is legitimate because
the Question's hypothesis is itself ``on a cone''; a reader who insists on $\Sin$-
definability strictly in $L$ with no oracle should note that the conclusion then holds
on the same cone on which the hypothesis is assumed, which is all the contradiction in
Corollary~\ref{cor:answer} requires.
\end{enumerate}
Subject to these standard relativisations, the answer to Question~1.2 is definitively
\textbf{No}.

\end{document}
